% Previous related work preserved (2026-09-26).
% \section{Related Work}
% 
% \textbf{Text-driven Human Motion Synthesize.} Text-to-Motion(T2M) generation aims to syhthesize 3D human motion from natural language. It has been dominated by two main paradigms. Diffusion-based methods such as MotionDiffuse~\citep{zhang2024motiondiffuse} and MDM~\citep{tevet2023mdm} generate motions by iterative denoising process with a text prompt. Following work such as ReMoDiffuse~\citep{zhang2023remodiffuse} and MLD~\citep{chen2023mld} improves the motion fidelity via retrieve argumentation. Meanwhile, autoregressive methods discretize motion into tokens and generate them sequentially. T2M-GPT~\citep{zhang2023t2mgpt} generates discrete motion tokens autoregressively, while MotionGPT~\citep{jiang2023motiongpt} further connects motion tokens to large language models. More recently, masked generative modeling has emerged as a strong alternative. MoMask~\citep{guo2024momask} employs residual vector quantization to represent motion as multi-layer discrete tokens. BAMM~\citep{pinyoanuntapong2024bamm} combines masked modeling with autoregressive generation, while MaskControl~\citep{pinyoanuntapong2025maskcontrol} introduces spatial control into masked motion synthesis. Recent works scale T2M models with larger architectures and datasets. HY-Motion~\citep{hymotion2025} and kimodo~\citep{Kimodo2026} scale DiT-based flow matching to the billion parameters and leverages large-scale motion pretraining to improve the motion quality. However, these methods rely primarily on text prompts or manipulating pose keyframes for control, limiting their ability to generate scene-aware human motions.
% 
% \textbf{Image-based Human Pose Regression.} Recovering 3D human pose and shape from images is commonly formulated as human mesh recovery (HMR). Most HMR methods regress the parameters of a statistical body model such as SMPL~\citep{SMPL:2015}. HMR~\citep{kanazawa2018hmr} uses iterative error feedback and adversarial pose priors to establish the end-to-end regression paradigm. SPIN~\citep{kolotouros2019spin} further improves HMR through optimization-in-the-loop training. CLIFF~\citep{li2022cliff} leverages the human location in the image to improve the mesh recovery. HybrIK~\citep{li2021hybrik} combines analytical inverse kinematics with neural networks for more accurate pose recovery. Humans in 4D~\citep{goel2023humansin4d} adopts a transformer-based architecture to reconstruct humans and track them over time. Recent HMR methods have expanded HMR methods toward whole-body~\citep{cai2023smpler} and multi-person reconstruction~\citep{baradel2024multihmr}. CameraHMR~\citep{patel2025camerahmr} improves camera modeling for accurate human placement in camera space. MetricHMSR~\citep{song2026metrichmsr} recovers metric human mesh and 3D scene geometry using perspective camera cues. OnlineHMR~\citep{zhao2026onlinehmr} further extends HMR to causal video processing, enabling online recovery of world-grounded human motion. The motions generated by these methods are highly dependent on the input images or videos, limiting their ability to produce diverse motions. With the rapid progress of video generation, another motion generation method use powerful video generation model such as Minimax-H3~\citep{minimaxh3} to synthesize diverse human videos and then extracted various motions by HMR-based methods~\citep{patel2025camerahmr}. However, two-stage pipeline requires considerable inference time.
% 
% \textbf{Scene-aware Human Pose Prediction.} Human motions are affected by the surrounding environments. POSA~\citep{hassan2021posa} leveraged human-scene contact and semantic relationships for plausible human placement in 3D scenes. SAMP~\citep{hassan2021samp} extends static motion modeling to dynamic motion generation. COINS~\citep{zhao2022coins} further introduces semantic control for compositional human-scene interactions. HUMANISE~\citep{wang2022humanise} further incorporates language conditions to generate semantically consistent human motions in 3D scenes. With the development of generative models, Diffusion-based approaches have been applied to human-object interaction synthesize. SceneDiffuser~\citep{jiang2024scenediffuser} conditions physical and orientation constraints to generate scene-aware human motions. OMOMO~\citep{li2023omomo} directly produces dynamic human-object interactions from object trajectories, without explicit specifications of the target interactions. Move as You Say~\citep{wang2024move} leverages scene affordance as an intermediate representation to generate language-guided human motions in 3D environments.TRUMANS~\citep{jiang2024trumans} models long-term human-scene interactions as a diffusion-based autoregressive model conditioned on scene context and actions. InfBaGel~\citep{zou2026infbagel} employs an iterative refinement strategy to generate interaction motions with dynamic scene awareness. Although these methods generate plausible interactions, they heavily rely on the 3D input, which limits their applicability. Move-in-2d~\citep{huang2025movein2d} adapts the motions to the visual environment to directly synthesize human motion just from a scene image and a text prompt. However, it lacks explicit control over the initial human pose, which makes it difficult to synthesize motions. 
% 
% \label{sec:related}

% Allow the introduction to flow into related work in the preprint.
\section{Related Work}
\label{sec:related}

% Before restoring masked-model details (2026-09-26).
% \noindent Text-conditioned human motion generation. Text-to-motion (T2M) methods generate 3D human motion from natural language. Diffusion-based methods such as MotionDiffuse~\citep{zhang2024motiondiffuse} and MDM~\citep{tevet2023mdm} denoise motion sequences, while MLD~\citep{chen2023mld} performs diffusion in a learned motion latent space and ReMoDiffuse~\citep{zhang2023remodiffuse} incorporates retrieved motion examples. Autoregressive methods, including T2M-GPT~\citep{zhang2023t2mgpt} and MotionGPT~\citep{jiang2023motiongpt}, generate discrete motion tokens sequentially. Masked models such as MoMask~\citep{guo2024momask}, BAMM~\citep{pinyoanuntapong2024bamm}, and MaskControl~\citep{pinyoanuntapong2025maskcontrol} further explore token prediction for motion generation and control. Recent work has also studied larger training datasets and more flexible motion constraints. HY-Motion~\citep{hymotion2025} scales DiT-based flow matching to a billion parameters, while Kimodo~\citep{Kimodo2026} trains a motion diffusion model on large-scale optical motion capture data and supports pose and trajectory constraints. These methods improve motion generation and control, but do not directly use scene images to model spatial relationships between a person and the surrounding environment.
\noindent \textbf{Text-conditioned human motion generation}. Text-to-motion (T2M) methods generate 3D human motion from natural language. Diffusion-based methods such as MotionDiffuse~\citep{zhang2024motiondiffuse} and MDM~\citep{tevet2023mdm} denoise motion sequences, while MLD~\citep{chen2023mld} performs diffusion in a learned motion latent space and ReMoDiffuse~\citep{zhang2023remodiffuse} incorporates retrieved motion examples. Autoregressive methods, including T2M-GPT~\citep{zhang2023t2mgpt} and MotionGPT~\citep{jiang2023motiongpt}, generate discrete motion tokens sequentially. Masked models such as MoMask~\citep{guo2024momask}, BAMM~\citep{pinyoanuntapong2024bamm}, and MaskControl~\citep{pinyoanuntapong2025maskcontrol} further explore token prediction for motion generation and control. MoMask uses residual vector quantization to represent motion with multiple token layers, while MaskControl incorporates spatial constraints. Recent work has also studied larger training datasets and more flexible motion constraints. HY-Motion~\citep{hymotion2025} scales DiT-based flow matching to a billion parameters, while Kimodo~\citep{Kimodo2026} trains a motion diffusion model on large-scale optical motion capture data and supports pose and trajectory constraints. These methods improve motion generation and control, but do not directly use scene images to model spatial relationships between a person and the surrounding environment.

% Before SceneDiffuser citation correction: \noindent Scene-conditioned human motion generation. Scene-conditioned methods incorporate environmental information into human pose and motion generation. Methods such as POSA~\citep{hassan2021posa}, SAMP~\citep{hassan2021samp}, COINS~\citep{zhao2022coins}, and HUMANISE~\citep{wang2022humanise} use 3D scene information to model human placement, contact, or motion. Other approaches explore different forms of interaction conditions: OMOMO~\citep{li2023omomo} uses object trajectories, Move as You Say~\citep{wang2024move} uses scene affordances, and TRUMANS~\citep{jiang2024trumans} conditions motion generation on scene context and actions. SceneDiffuser~\citep{jiang2024scenediffuser} and InfBaGel~\citep{zou2026infbagel} further study scene-conditioned generation and refinement for human interactions. These methods explore how scene information can guide motion beyond textual descriptions. Move-in-2D~\citep{huang2025movein2d} generates motion from a background image and text, and uses the generated motion to guide subsequent video generation. Our work leverages the priors learned during video world model pretraining to improve motion generalization across diverse environments.
% Before agreed claim-scope revision: \noindent Scene-conditioned human motion generation. Scene-conditioned methods incorporate environmental information into human pose and motion generation. Methods such as POSA~\citep{hassan2021posa}, SAMP~\citep{hassan2021samp}, COINS~\citep{zhao2022coins}, and HUMANISE~\citep{wang2022humanise} use 3D scene information to model human placement, contact, or motion. Other approaches explore different forms of interaction conditions: OMOMO~\citep{li2023omomo} uses object trajectories, Move as You Say~\citep{wang2024move} uses scene affordances, and TRUMANS~\citep{jiang2024trumans} conditions motion generation on scene context and actions. SceneDiffuser~\citep{huang2023scenediffuser} and InfBaGel~\citep{zou2026infbagel} further study scene-conditioned generation and refinement for human interactions. These methods explore how scene information can guide motion beyond textual descriptions. Move-in-2D~\citep{huang2025movein2d} generates motion from a background image and text, and uses the generated motion to guide subsequent video generation. Our work leverages the priors learned during video world model pretraining to improve motion generalization across diverse environments.
\noindent \textbf{Scene-conditioned human motion generation.} Scene-conditioned methods incorporate environmental information into human pose and motion generation. Methods such as POSA~\citep{hassan2021posa}, SAMP~\citep{hassan2021samp}, COINS~\citep{zhao2022coins}, and HUMANISE~\citep{wang2022humanise} use 3D scene information to model human placement, contact, or motion. Other approaches explore different forms of interaction conditions: OMOMO~\citep{li2023omomo} uses object trajectories, Move as You Say~\citep{wang2024move} uses scene affordances, and TRUMANS~\citep{jiang2024trumans} conditions motion generation on scene context and actions. SceneDiffuser~\citep{huang2023scenediffuser} and InfBaGel~\citep{zou2026infbagel} further study scene-conditioned generation and refinement for human interactions. These methods explore how scene information can guide motion beyond textual descriptions. Move-in-2D~\citep{huang2025movein2d} generates motion from a background image and text, and uses the generated motion to guide subsequent video generation. Our work studies the adaptation of a pretrained video world model for image-conditioned joint video and 3D motion generation.

\noindent \textbf{Human motion recovery from images and videos.} Human motion recovery estimates 3D body pose and shape from observed images or videos, often using a parametric body model such as SMPL~\citep{SMPL:2015}. HMR~\citep{kanazawa2018hmr}, SPIN~\citep{kolotouros2019spin}, CLIFF~\citep{li2022cliff}, and HybrIK~\citep{li2021hybrik} improve recovery through learned pose priors, optimization during training, image-location cues, and kinematic constraints. More recent methods study transformer-based recovery, whole-body and multi-person reconstruction, and camera and temporal modeling, including Humans in 4D~\citep{goel2023humansin4d}, SMPLer-X~\citep{cai2023smpler}, Multi-HMR~\citep{baradel2024multihmr}, CameraHMR~\citep{patel2025camerahmr}, MetricHMSR~\citep{song2026metrichmsr}, and OnlineHMR~\citep{zhao2026onlinehmr}. In our framework, CameraHMR provides motion estimates for real training videos and the initial pose from the input image. Motion recovery can also be applied to generated videos, forming a two-stage motion-generation pipeline, but the additional recovery stage increases inference cost. In contrast, World2Motion jointly generates motion and video, achieving approximately $3.3\times$ faster inference, better motion--text alignment, and lower motion jitter than the MiniMax-H3~\citep{minimaxh3} + CameraHMR~\citep{patel2025camerahmr} baseline.

\noindent \textbf{Video world models and joint video--motion generation.} Large-scale video models, such as Wan~\citep{wan2025wan} and MiniMax-H3~\citep{minimaxh3}, learn visual priors about human motion and environmental interactions. Cosmos 3~\citep{agarwal2026cosmos} extends multimodal generation to video and action within a Mixture-of-Transformers architecture. Its action modeling includes egocentric head and hand motion, and its action training inherits the vision noise schedule. This provides a starting point for adapting video world models to full-body 3D motion generation. Related work has also explored joint video--motion generation. CoMoVi~\citep{zhao2026comovi} uses separate diffusion branches for RGB video and 2D motion representations, together with 3D--2D cross-attention to predict 3D motion within the same denoising loop. World2Motion adapts the existing multimodal backbone of Cosmos 3 for joint video--motion generation. We combine synthetic and real video--motion training pairs and propose a shift-decoupled noise schedule that assigns modality-specific noise levels through shared denoising progress to improve temporal stability.

% Keep section transitions continuous in the preprint.
% Let floats share the following text page.